
% ════════════════════════════════════════════
% 模型准备（5.2.2.模型准备.tex，三级子块）—— 写作提示词见本文件末尾【本章节写作提示词】。
% ════════════════════════════════════════════
\subsubsection{模型准备}

\textbf{1.数据预处理}

本问沿用问题一的可观测性判定框架，将告警 T2 按式(1)换算为坐标 $(11.03,72.5)$、高度 35 m。为便于站对枚举与时序仿真，首先对区域进行 1 m 网格离散，提取同时满足三项约束的可观测点集，并记录每点的目标方位角与水平距离，形成候选观测站表。由于目标 T2 位于建筑 B3 正上方，其可观测区域受 B3 阴影影响主要分布在西侧与北侧，东侧与南侧需满足更大的距离下界，这一结构将直接决定最优站对的方位配置。

\textbf{2.算法选择}

针对本问最优站对与并发调度，需在穷举搜索、贪心启发式与数学规划法之间择优，从求解质量、计算复杂度与可解释性三个维度对比其适用性，如表\ref{tab:algo2}所示。

\begin{longtable}{>{\centering\arraybackslash}p{0.20\textwidth}>{\centering\arraybackslash}p{0.22\textwidth}>{\centering\arraybackslash}p{0.22\textwidth}>{\centering\arraybackslash}p{0.24\textwidth}}
  \caption{问题二候选算法对比}
  \label{tab:algo2} \\
  \toprule
  算法 & 求解质量 & 计算复杂度 & 可解释性 \\
  \midrule
  \endfirsthead
  \caption{问题二候选算法对比（续）} \\
  \toprule
  算法 & 求解质量 & 计算复杂度 & 可解释性 \\
  \midrule
  \endhead
  \bottomrule
  \endfoot
  穷举搜索 & 全局最优 & 小规模可接受 & 强 \\
  贪心启发式 & 局部最优 & 低 & 强 \\
  数学规划法 & 全局最优 & 中 & 中 \\
\end{longtable}

如表\ref{tab:algo2}所示，本问任务数为 6、车辆数为 5，属小规模组合问题，穷举搜索可在秒级内求得全局最优，且避免了机器学习类方法在小样本下易过拟合的问题\cite{zhou2016jiqi}。综合以上分析，本问选择穷举搜索作为调度主算法，并以独立定义的最近指派启发式作为基准对比。
